% Options for packages loaded elsewhere
\PassOptionsToPackage{unicode}{hyperref}
\PassOptionsToPackage{hyphens}{url}
\documentclass[
]{article}
\usepackage{xcolor}
\usepackage[letterpaper,top=2cm,bottom=2cm,left=1.5cm,right=1.5cm]{geometry}
\usepackage{amsmath,amssymb}
\setcounter{secnumdepth}{-\maxdimen} % remove section numbering
\usepackage{iftex}
\ifPDFTeX
  \usepackage[T1]{fontenc}
  \usepackage[utf8]{inputenc}
  \usepackage{textcomp} % provide euro and other symbols
\else % if luatex or xetex
  \usepackage{unicode-math} % this also loads fontspec
  \defaultfontfeatures{Scale=MatchLowercase}
  \defaultfontfeatures[\rmfamily]{Ligatures=TeX,Scale=1}
\fi
\usepackage{lmodern}
\ifPDFTeX\else
  % xetex/luatex font selection
\fi
% Use upquote if available, for straight quotes in verbatim environments
\IfFileExists{upquote.sty}{\usepackage{upquote}}{}
\IfFileExists{microtype.sty}{% use microtype if available
  \usepackage[]{microtype}
  \UseMicrotypeSet[protrusion]{basicmath} % disable protrusion for tt fonts
}{}
\makeatletter
\@ifundefined{KOMAClassName}{% if non-KOMA class
  \IfFileExists{parskip.sty}{%
    \usepackage{parskip}
  }{% else
    \setlength{\parindent}{0pt}
    \setlength{\parskip}{6pt plus 2pt minus 1pt}}
}{% if KOMA class
  \KOMAoptions{parskip=half}}
\makeatother
\usepackage{color}
\usepackage{fancyvrb}
\newcommand{\VerbBar}{|}
\newcommand{\VERB}{\Verb[commandchars=\\\{\}]}
\DefineVerbatimEnvironment{Highlighting}{Verbatim}{commandchars=\\\{\}}
% Add ',fontsize=\small' for more characters per line
\newenvironment{Shaded}{}{}
\newcommand{\AlertTok}[1]{\textcolor[rgb]{1.00,0.00,0.00}{\textbf{#1}}}
\newcommand{\AnnotationTok}[1]{\textcolor[rgb]{0.38,0.63,0.69}{\textbf{\textit{#1}}}}
\newcommand{\AttributeTok}[1]{\textcolor[rgb]{0.49,0.56,0.16}{#1}}
\newcommand{\BaseNTok}[1]{\textcolor[rgb]{0.25,0.63,0.44}{#1}}
\newcommand{\BuiltInTok}[1]{\textcolor[rgb]{0.00,0.50,0.00}{#1}}
\newcommand{\CharTok}[1]{\textcolor[rgb]{0.25,0.44,0.63}{#1}}
\newcommand{\CommentTok}[1]{\textcolor[rgb]{0.38,0.63,0.69}{\textit{#1}}}
\newcommand{\CommentVarTok}[1]{\textcolor[rgb]{0.38,0.63,0.69}{\textbf{\textit{#1}}}}
\newcommand{\ConstantTok}[1]{\textcolor[rgb]{0.53,0.00,0.00}{#1}}
\newcommand{\ControlFlowTok}[1]{\textcolor[rgb]{0.00,0.44,0.13}{\textbf{#1}}}
\newcommand{\DataTypeTok}[1]{\textcolor[rgb]{0.56,0.13,0.00}{#1}}
\newcommand{\DecValTok}[1]{\textcolor[rgb]{0.25,0.63,0.44}{#1}}
\newcommand{\DocumentationTok}[1]{\textcolor[rgb]{0.73,0.13,0.13}{\textit{#1}}}
\newcommand{\ErrorTok}[1]{\textcolor[rgb]{1.00,0.00,0.00}{\textbf{#1}}}
\newcommand{\ExtensionTok}[1]{#1}
\newcommand{\FloatTok}[1]{\textcolor[rgb]{0.25,0.63,0.44}{#1}}
\newcommand{\FunctionTok}[1]{\textcolor[rgb]{0.02,0.16,0.49}{#1}}
\newcommand{\ImportTok}[1]{\textcolor[rgb]{0.00,0.50,0.00}{\textbf{#1}}}
\newcommand{\InformationTok}[1]{\textcolor[rgb]{0.38,0.63,0.69}{\textbf{\textit{#1}}}}
\newcommand{\KeywordTok}[1]{\textcolor[rgb]{0.00,0.44,0.13}{\textbf{#1}}}
\newcommand{\NormalTok}[1]{#1}
\newcommand{\OperatorTok}[1]{\textcolor[rgb]{0.40,0.40,0.40}{#1}}
\newcommand{\OtherTok}[1]{\textcolor[rgb]{0.00,0.44,0.13}{#1}}
\newcommand{\PreprocessorTok}[1]{\textcolor[rgb]{0.74,0.48,0.00}{#1}}
\newcommand{\RegionMarkerTok}[1]{#1}
\newcommand{\SpecialCharTok}[1]{\textcolor[rgb]{0.25,0.44,0.63}{#1}}
\newcommand{\SpecialStringTok}[1]{\textcolor[rgb]{0.73,0.40,0.53}{#1}}
\newcommand{\StringTok}[1]{\textcolor[rgb]{0.25,0.44,0.63}{#1}}
\newcommand{\VariableTok}[1]{\textcolor[rgb]{0.10,0.09,0.49}{#1}}
\newcommand{\VerbatimStringTok}[1]{\textcolor[rgb]{0.25,0.44,0.63}{#1}}
\newcommand{\WarningTok}[1]{\textcolor[rgb]{0.38,0.63,0.69}{\textbf{\textit{#1}}}}
\usepackage{longtable,booktabs,array}
\usepackage{caption}
\captionsetup[table]{skip=6pt}
\newcounter{none} % for unnumbered tables
\usepackage{calc} % for calculating minipage widths
% Correct order of tables after \paragraph or \subparagraph
\usepackage{etoolbox}
\makeatletter
\patchcmd\longtable{\par}{\if@noskipsec\mbox{}\fi\par}{}{}
\makeatother
% Allow footnotes in longtable head/foot
\IfFileExists{footnotehyper.sty}{\usepackage{footnotehyper}}{\usepackage{footnote}}
\makesavenoteenv{longtable}
\setlength{\emergencystretch}{3em} % prevent overfull lines
\providecommand{\tightlist}{%
  \setlength{\itemsep}{0pt}\setlength{\parskip}{0pt}}
\usepackage{bookmark}
\IfFileExists{xurl.sty}{\usepackage{xurl}}{} % add URL line breaks if available
\urlstyle{same}
% fallback for those not using the hyperref driver hyperxmp:
\makeatletter
\@ifundefined{xmpquote}{\newcommand{\xmpquote}[1]{#1}}{}
\makeatother
\hypersetup{
  pdftitle={Regression Analysis - Lecture Notes},
  pdfauthor={MA2003B},
  hidelinks,
  pdfcreator={LaTeX via pandoc}}

\title{Regression Analysis - Lecture Notes}
\author{MA2003B}
\date{}

\begin{document}
\maketitle

\textbf{Regression Analysis}

Comprehensive Study Guide \& Theoretical Notes

MA2003B --- Application of Multivariate Methods in Data Science

Tecnológico de Monterrey

\begin{center}\rule{0.5\linewidth}{0.5pt}\end{center}

\textbf{ }

\section{Course and Chapter Overview}

\textbf{ }

\subsection{Course Context}

In the \textbf{MA2003B} curriculum, regression analysis serves as the
essential stepping stone connecting univariate inference to
high-dimensional multivariate models. This chapter develops both
classical Ordinary Least Squares (OLS) foundation and modern diagnostics
for building predictive statistical relationships.

\textbf{ }

\subsection{Learning Objectives}

By the end of this module, students will be able to:

\begin{itemize}
\item
  Derive and interpret Simple and Multiple Linear Regression models.
\item
  Decompose sums of squares via ANOVA tables and evaluate overall and
  partial \(F\)-tests.
\item
  Construct and contrast confidence intervals for mean response vs.
  prediction intervals for new observations.
\item
  Conduct residual diagnostic procedures to evaluate normality,
  linearity, and homoskedasticity.
\item
  Implement forward selection, backward elimination, and
  information-criterion-based model searches (AIC, BIC).
\item
  Incorporate non-linear polynomial terms and logarithmic
  transformations.
\item
  Identify multi-collinearity using Variance Inflation Factors (VIF).
\item
  Detect high-leverage points and influential outliers using Cook's
  Distance and Studentized residuals.
\item
  Address heteroskedasticity through Weighted Least Squares (WLS) and
  robust standard errors (HC3).
\end{itemize}

---

\textbf{ }

\section{1.1 Simple Linear Regression}

\textbf{ }

\subsection{Mathematical Formulation}

Simple Linear Regression models the linear association between a
continuous scalar response variable \(Y\) and a single explanatory
predictor variable \(X\):

\[y_{i} = \beta_{0} + \beta_{1}x_{i} + \varepsilon_{i},\quad i = 1,2,\ldots,n\]

where:

\begin{itemize}
\item
  \(\beta_{0} \in \mathbb{R}\) is the population intercept parameter.
\item
  \(\beta_{1} \in \mathbb{R}\) is the population slope coefficient
  (expected marginal change in \(Y\) per unit change in \(X\)).
\item
  \(\varepsilon_{i}\) is the unobserved random error term satisfying the
  classical Gauss-Markov assumptions: \[\begin{array}{r}
  \mathbb{E}\left\lbrack \varepsilon_{i} \right\rbrack = 0,\quad\operatorname{Var}(\varepsilon_{i}) = \sigma^{2},\quad\operatorname{Cov}(\varepsilon_{i},\varepsilon_{j}) = 0 \\
  (i \neq j)
  \end{array}\]
\end{itemize}

Under the normality assumption,
\(\varepsilon_{i} \sim^{\text{iid}}\mathcal{N}(0,\sigma^{2})\).

\textbf{ }

\subsection{Ordinary Least Squares (OLS) Derivation}

The OLS estimators
\(\left( {\hat{\beta}}_{0},{\hat{\beta}}_{1} \right)\) minimize the
Residual Sum of Squares:

\[S\left( \beta_{0},\beta_{1} \right) = \sum_{i = 1}^{n}\left( y_{i} - \beta_{0} - \beta_{1}x_{i} \right)^{2}\]

Setting partial derivatives to zero yields the normal equations:
\[\frac{\partial S}{\partial\beta_{0}} = - 2\sum_{i = 1}^{n}\left( y_{i} - \beta_{0} - \beta_{1}x_{i} \right) = 0 \Longrightarrow {\hat{\beta}}_{0} = \overline{y} - {\hat{\beta}}_{1}\overline{x}\]
\[\frac{\partial S}{\partial\beta_{1}} = - 2\sum_{i = 1}^{n}x_{i}\left( y_{i} - \beta_{0} - \beta_{1}x_{i} \right) = 0 \Longrightarrow {\hat{\beta}}_{1} = \frac{\sum_{i = 1}^{n}\left( x_{i} - \overline{x} \right)\left( y_{i} - \overline{y} \right)}{\sum_{i = 1}^{n}\left( x_{i} - \overline{x} \right)^{2}} = \frac{S_{xy}}{S_{xx}}\]

\textbf{Info: } \textbf{Gauss-Markov Theorem:} Under the assumptions of
linearity, strict exogeneity, homoskedasticity, and no autocorrelation,
the OLS estimator \(\hat{\mathbf{\beta}}\) is the \textbf{Best Linear
Unbiased Estimator (BLUE)}.

---

\textbf{ }

\section{1.2 ANOVA and Confidence Intervals for Predictions}

\textbf{ }

\subsection{Sum of Squares Decomposition}

The total sample variation in \(Y\) is partitioned into explained
(regression) and unexplained (residual) components:

\[\underset{\text{ SST (Total)}}{\underbrace{\sum_{i = 1}^{n}\left( y_{i} - \overline{y} \right)^{2}}} = \underset{\text{ SSR (Regression)}}{\underbrace{\sum_{i = 1}^{n}\left( {\hat{y}}_{i} - \overline{y} \right)^{2}}} + \underset{\text{ SSE (Error)}}{\underbrace{\sum_{i = 1}^{n}\left( y_{i} - {\hat{y}}_{i} \right)^{2}}}\]

Degrees of freedom: \(\text{df}_{\text{Total}} = n - 1\),
\(\text{df}_{\text{Regression}} = 1\) (or \(k\) predictors),
\(\text{df}_{\text{Error}} = n - k - 1\).

{\def\LTcaptype{none} % do not increment counter
\begin{longtable}[]{@{}
  >{\centering\arraybackslash}p{(\linewidth - 8\tabcolsep) * \real{0.2500}}
  >{\centering\arraybackslash}p{(\linewidth - 8\tabcolsep) * \real{0.1875}}
  >{\centering\arraybackslash}p{(\linewidth - 8\tabcolsep) * \real{0.1875}}
  >{\centering\arraybackslash}p{(\linewidth - 8\tabcolsep) * \real{0.1875}}
  >{\centering\arraybackslash}p{(\linewidth - 8\tabcolsep) * \real{0.1875}}@{}}
\toprule\noalign{}
\begin{minipage}[b]{\linewidth}\centering
\textbf{Source}
\end{minipage} & \begin{minipage}[b]{\linewidth}\centering
\textbf{Sum of Squares}
\end{minipage} & \begin{minipage}[b]{\linewidth}\centering
\textbf{Degrees of Freedom}
\end{minipage} & \begin{minipage}[b]{\linewidth}\centering
\textbf{Mean Square (MS)}
\end{minipage} & \begin{minipage}[b]{\linewidth}\centering
\textbf{\(F\)-Statistic}
\end{minipage} \\
\midrule\noalign{}
\endhead
\bottomrule\noalign{}
\endlastfoot
Regression & SSR & \(k\) & \(\text{MSR } = \frac{\text{ SSR}}{k}\) &
\(F = \frac{\text{MSR}}{\text{ MSE}}\) \\
Residual (Error) & SSE & \(n - k - 1\) &
\(\text{MSE } = \frac{\text{ SSE}}{n - k - 1}\) & \\
Total & SST & \(n - 1\) & & \\
\end{longtable}
}

Coefficient of determination:
\[R^{2} = \frac{\text{SSR}}{\text{ SST}} = 1 - \frac{\text{SSE}}{\text{ SST}}\]

\textbf{ }

\subsection{Mean Response Confidence Interval vs. Individual Prediction
Interval}

For a specified predictor value \(x_{0}\):

\textbf{1. Confidence Interval for the Expected Mean
\(\mathbb{E}\left\lbrack Y~|~X = x_{0} \right\rbrack\):}
\[{\hat{y}}_{0} \pm t_{n - 2,1 - \frac{\alpha}{2}} \cdot s\sqrt{\frac{1}{n} + \frac{\left( x_{0} - \overline{x} \right)^{2}}{S_{xx}}}\]

\textbf{2. Prediction Interval for a Single New Observation
\(Y_{\text{new}}\):}
\[{\hat{y}}_{0} \pm t_{n - 2,1 - \frac{\alpha}{2}} \cdot s\sqrt{1 + \frac{1}{n} + \frac{\left( x_{0} - \overline{x} \right)^{2}}{S_{xx}}}\]

\textbf{Tip: } Prediction intervals are always strictly wider than
confidence intervals because they account for both model parameter
uncertainty and individual observation variance \(\sigma^{2}\).

---

\textbf{ }

\section{1.3 Residual Analysis \& Condition of Normality}

\textbf{ }

\subsection{Residual Definitions}

\begin{itemize}
\item
  \textbf{Raw Residuals:} \(e_{i} = y_{i} - {\hat{y}}_{i}\).
\item
  \textbf{Standardized Residuals:} \(d_{i} = \frac{e_{i}}{s}\).
\item
  \textbf{Studentized (Internally Studentized) Residuals:}
  \[r_{i} = \frac{e_{i}}{s\sqrt{1 - h_{ii}}}\] where
  \(h_{ii} = \mathbf{x}_{i}^{T}\left( \mathbf{X}^{T}\mathbf{X} \right)^{- 1}\mathbf{x}_{i}\)
  is the \(i\)-th diagonal element of the hat matrix
  \(\mathbf{H} = {\mathbf{X}(\mathbf{X}^{T}\mathbf{X})}^{- 1}\mathbf{X}^{T}\).
\end{itemize}

\textbf{ }

\subsection{Diagnostic Tools}

\begin{itemize}
\item
  \textbf{Residuals vs. Fitted Values Plot:} Evaluates linearity (should
  show horizontal band around zero) and homoskedasticity (constant
  vertical spread).
\item
  \textbf{Normal Q-Q Plot:} Plots sample quantiles of studentized
  residuals against theoretical normal quantiles. Departures from the
  \(45^{\circ}\) reference line indicate skewness or heavy tails.
\item
  \textbf{Shapiro-Wilk Test:} Statistical hypothesis test with null
  hypothesis \(H_{0}\): residuals are normally distributed.
\end{itemize}

---

\textbf{ }

\section{1.4 \& 1.5 Variable Selection Methods}

\textbf{ }

\subsection{Forward Selection}

\begin{enumerate}
\item
  Start with the intercept-only null model: \(y = \beta_{0}\).
\item
  For each candidate predictor \(X_{j}\), compute the \(F\)-to-enter
  statistic (or \(p\)-value) when added to the current model.
\item
  Add the predictor with the lowest \(p\)-value if
  \(p < \alpha_{\text{enter}}\) (typically \(0.05\)).
\item
  Repeat until no remaining predictor meets the entry criterion.
\end{enumerate}

\textbf{ }

\subsection{Backward Elimination}

\begin{enumerate}
\item
  Start with the full model containing all \(p\) candidate predictors.
\item
  For each predictor in the model, calculate the partial \(F\)-test or
  \(t\)-test \(p\)-value.
\item
  Remove the variable with the highest \(p\)-value if
  \(p > \alpha_{\text{remove}}\) (typically \(0.10\)).
\item
  Re-fit and repeat until all remaining variables are statistically
  significant.
\end{enumerate}

\textbf{ }

\subsection{Information Criteria Comparison}

\begin{itemize}
\item
  \textbf{Akaike Information Criterion (AIC):}
  \(\text{AIC } = 2k - 2\ln(\hat{L}) \approx n\ln(\frac{\text{SSE}}{n}) + 2k\).
  Penalizes model complexity.
\item
  \textbf{Bayesian Information Criterion (BIC):}
  \(\text{BIC } = k\ln(n) - 2\ln(\hat{L}) \approx n\ln(\frac{\text{SSE}}{n}) + k\ln(n)\).
  Imposes a stronger penalty for large \(n\), favoring more parsimonious
  models.
\end{itemize}

---

\textbf{ }

\section{1.6 Non-Linear Regression \& Transformations}

When the true relationship between \(Y\) and \(X\) exhibits curvature:

\textbf{ }

\subsection{Polynomial Models}

\[y_{i} = \beta_{0} + \beta_{1}x_{i} + \beta_{2}x_{i}^{2} + \ldots + \beta_{d}x_{i}^{d} + \varepsilon_{i}\]
Linear in parameters \(\beta\), estimated using standard OLS.

\textbf{ }

\subsection{Logarithmic Transformations}

The percentage-change interpretations below are \textbf{local
(small-coefficient) approximations}, valid for \(\beta_{1}\) near 0 or
for describing an infinitesimal change in \(X\). For the exact effect of
a full one-unit (or 1\%) change, use the exact formula alongside each
approximation.

\begin{itemize}
\item
  \textbf{Log-Linear (Exponential growth):}
  \(\ln(y_{i}) = \beta_{0} + \beta_{1}x_{i} + \varepsilon_{i}\).
  Approximate: \(100 \cdot \beta_{1}\%\) change in \(Y\) per unit \(X\).
  Exact one-unit effect: \(100\left( e^{\beta_{1}} - 1 \right)\%\).
\item
  \textbf{Linear-Log:}
  \(y_{i} = \beta_{0} + \beta_{1}\ln(x_{i}) + \varepsilon_{i}\).
  Approximate: \(\frac{\beta_{1}}{100}\) change in \(Y\) per \(1\%\)
  increase in \(X\) (exact for small \(\Delta X/X\); for a full 1\%
  increase, \(\Delta Y \approx \beta_{1}\ln(1.01)\)).
\item
  \textbf{Log-Log (Elasticity):}
  \(\ln(y_{i}) = \beta_{0} + \beta_{1}\ln(x_{i}) + \varepsilon_{i}\).
  Approximate: \(\beta_{1}\%\) change in \(Y\) per \(1\%\) increase in
  \(X\). Exact one-percent effect:
  \(100\left( 1.01^{\beta_{1}} - 1 \right)\%\). \(\beta_{1}\) is the
  \textbf{elasticity} of \(Y\) with respect to \(X\) exactly, regardless
  of the size of the change.
\end{itemize}

---

\textbf{ }

\section{1.7 Multiple Linear Regression}

\textbf{ }

\subsection{Matrix Formulation}

\[\mathbf{y} = \mathbf{X}\mathbf{\beta} + \mathbf{\epsilon}\]

where \(\mathbf{y} \in \mathbb{R}^{n \times 1}\),
\(\mathbf{X} \in \mathbb{R}^{n \times (k + 1)}\),
\(\mathbf{\beta} \in \mathbb{R}^{(k + 1) \times 1}\),
\(\mathbf{\epsilon} \sim \mathcal{N}_{n\left( \mathbf{0},\sigma^{2}\mathbf{I}_{n} \right)}\).

\textbf{ }

\subsection{Parameter Estimation \& Covariance Matrix}

\[\hat{\mathbf{\beta}} = \left( \mathbf{X}^{T}\mathbf{X} \right)^{- 1}\mathbf{X}^{T}\mathbf{y}\]
\[\operatorname{Cov}(\hat{\mathbf{\beta}}) = \sigma^{2}\left( \mathbf{X}^{T}\mathbf{X} \right)^{- 1}\quad\left( \text{Sample estimator:  }s^{2}\left( \mathbf{X}^{T}\mathbf{X} \right)^{- 1} \right)\]

\textbf{ }

\subsection{Multicollinearity \& Variance Inflation Factor (VIF)}

Multicollinearity occurs when two or more predictor variables in
\(\mathbf{X}\) are highly linearly correlated. For predictor \(X_{j}\):
\[\text{ VIF}_{j} = \frac{1}{1 - R_{j}^{2}}\] where \(R_{j}^{2}\) is the
coefficient of determination from regressing \(X_{j}\) on all other
remaining predictors.

\begin{itemize}
\item
  \(\text{VIF } = 1\): No collinearity.
\item
  \(\text{VIF } > 5\): Moderate collinearity requiring attention.
\item
  \(\text{VIF } > 10\): Severe collinearity causing unstable coefficient
  estimates and inflated standard errors.
\end{itemize}

---

\textbf{ }

\section{1.8 Aberrant Data \& Heteroskedasticity Problems}

\textbf{ }

\subsection{Leverage and Influence Diagnostics}

\begin{itemize}
\item
  \textbf{Hat Matrix Leverage (\(h_{ii}\)):} Measures how far
  observation \(i\)`s predictors are from the multivariate centroid:
  \[\overline{h} = \frac{k + 1}{n}\quad\left( \text{High leverage if  }h_{ii} > 2\overline{h} \right)\]
\item
  \textbf{Cook's Distance (\(D_{i}\)):} Measures the aggregate shift in
  fitted values when observation \(i\) is deleted:
  \[D_{i} = \frac{\sum_{j = 1}^{n}\left( {\hat{y}}_{j} - {\hat{y}}_{j(i)} \right)^{2}}{(k + 1)s^{2}} = \frac{r_{i}^{2}}{k + 1}\left( \frac{h_{ii}}{1 - h_{ii}} \right)\]
  Values of \(D_{i} > 1.0\) (or \(> \frac{4}{n}\)) indicate highly
  influential points.
\end{itemize}

\textbf{ }

\subsection{Heteroskedasticity Detection \& Remedies}

\begin{itemize}
\item
  \textbf{Breusch-Pagan Test:} Regresses squared standardized residuals
  \(\frac{e_{i}^{2}}{{\hat{\sigma}}^{2}}\) on predictors \(\mathbf{X}\).
  Under \(H_{0}\) (homoskedasticity), the test statistic follows
  \(\chi_{k}^{2}\).
\item
  \textbf{Weighted Least Squares (WLS):} When
  \(\operatorname{Var}(\varepsilon_{i}) = \sigma^{2}w_{i}^{- 1}\),
  minimize
  \(\sum w_{i}\left( y_{i} - \mathbf{x}_{i}^{T}\mathbf{\beta} \right)^{2}\):
  \[{\hat{\mathbf{\beta}}}_{\text{WLS}} = \left( \mathbf{X}^{T}\mathbf{W}\mathbf{X} \right)^{- 1}\mathbf{X}^{T}\mathbf{W}\mathbf{y},\quad\mathbf{W} = \operatorname{diag}(w_{1},\ldots,w_{n})\]
\item
  \textbf{Heteroskedasticity-Consistent Covariance (White / HC3):}
  \[\operatorname{Cov}_{HC3}\left( \hat{\mathbf{\beta}} \right) = \left( \mathbf{X}^{T}\mathbf{X} \right)^{- 1}\left( \sum_{i = 1}^{n}\frac{e_{i}^{2}}{\left( 1 - h_{ii} \right)^{2}}\mathbf{x}_{i}\mathbf{x}_{i}^{T} \right)\left( \mathbf{X}^{T}\mathbf{X} \right)^{- 1}\]
\end{itemize}

---

\textbf{ }

\section{Practical Python Implementation}

\begin{Shaded}
\begin{Highlighting}[]
\ImportTok{import}\NormalTok{ statsmodels.api }\ImportTok{as}\NormalTok{ sm}
\ImportTok{import}\NormalTok{ statsmodels.formula.api }\ImportTok{as}\NormalTok{ smf}
\ImportTok{from}\NormalTok{ statsmodels.stats.outliers\_influence }\ImportTok{import}\NormalTok{ variance\_inflation\_factor, OLSInfluence}
\ImportTok{from}\NormalTok{ statsmodels.stats.diagnostic }\ImportTok{import}\NormalTok{ het\_breuschpagan}
\ImportTok{import}\NormalTok{ pandas }\ImportTok{as}\NormalTok{ pd}
\ImportTok{import}\NormalTok{ numpy }\ImportTok{as}\NormalTok{ np}

\CommentTok{\# 1. Fit Multiple Linear Regression}
\NormalTok{model }\OperatorTok{=}\NormalTok{ smf.ols(}\StringTok{"sale\_price \textasciitilde{} sqft\_living + bedrooms + bathrooms + house\_age\_years + np.power(house\_age\_years, 2) + dist\_city\_center\_km + building\_grade"}\NormalTok{, data}\OperatorTok{=}\NormalTok{df).fit()}
\BuiltInTok{print}\NormalTok{(model.summary())}

\CommentTok{\# 2. Multicollinearity Assessment (VIF)}
\NormalTok{X }\OperatorTok{=}\NormalTok{ model.model.exog}
\NormalTok{vif\_data }\OperatorTok{=}\NormalTok{ pd.DataFrame(\{}
    \StringTok{"Variable"}\NormalTok{: model.model.exog\_names,}
    \StringTok{"VIF"}\NormalTok{: [variance\_inflation\_factor(X, i) }\ControlFlowTok{for}\NormalTok{ i }\KeywordTok{in} \BuiltInTok{range}\NormalTok{(X.shape[}\DecValTok{1}\NormalTok{])]}
\NormalTok{\})}
\BuiltInTok{print}\NormalTok{(vif\_data)}

\CommentTok{\# 3. Breusch{-}Pagan Test for Heteroskedasticity}
\NormalTok{bp\_test }\OperatorTok{=}\NormalTok{ het\_breuschpagan(model.resid, model.model.exog)}
\BuiltInTok{print}\NormalTok{(}\SpecialStringTok{f"BP Lagrange Multiplier p{-}value: }\SpecialCharTok{\{}\NormalTok{bp\_test[}\DecValTok{1}\NormalTok{]}\SpecialCharTok{:.4e\}}\SpecialStringTok{"}\NormalTok{)}

\CommentTok{\# 4. Robust Standard Errors (HC3)}
\NormalTok{robust\_model }\OperatorTok{=}\NormalTok{ model.get\_robustcov\_results(cov\_type}\OperatorTok{=}\StringTok{"HC3"}\NormalTok{)}
\BuiltInTok{print}\NormalTok{(robust\_model.summary())}
\end{Highlighting}
\end{Shaded}


\end{document}

